\providecommand{\tightlist}{%
  \setlength{\itemsep}{0pt}\setlength{\parskip}{0pt}}
\documentclass{article}
\usepackage{arxiv}
\usepackage{amsmath}
\usepackage{graphicx}
\usepackage{hyperref}
\usepackage{fontspec}
\usepackage{unicode-math}
\setmainfont{Linux Libertine O}
\setmathfont{XITS Math}
\usepackage{xcolor}
\usepackage{booktabs}
\usepackage{multirow}

\title{Compiled Kernels Amplify Data Ordering Effects\\in Language Model Pretraining}
\author{Augustin Chan\\
\texttt{aug@iterative.day}\\
Independent Researcher}
\date{\today}

\begin{document}
\maketitle

\begin{abstract}
We report an interaction between \texttt{torch.compile} and training data ordering
that produces curriculum-like effects absent on non-compiled frameworks.
In controlled experiments training a 3M-parameter GPT model on identical data,
random shuffling of a 64-batch buffer improves validation bits-per-byte by
0.106 on CUDA with \texttt{torch.compile} --- 2.7$\times$ the seed noise floor ---
while the same shuffling produces improvements of 0.008--0.037 on Apple MLX,
within measurement noise.
The effect persists across two independent difficulty metrics (token diversity
and compression ratio), ruling out metric artifacts.
We hypothesize that \texttt{torch.compile}'s kernel fusion and graph optimization
specialize execution paths to sequential data access patterns created by
best-fit document packing; shuffling disrupts these patterns, forcing
more general computation.
To our knowledge, this interaction between framework compilation and
data ordering has not been previously documented.
Our results have practical implications: curriculum learning experiments
conducted under \texttt{torch.compile} may report spurious effects that
do not generalize to other frameworks or deployment settings.
\end{abstract}

\textbf{Keywords:} torch.compile, curriculum learning, data ordering,
compilation artifacts, language model pretraining, MLX, kernel optimization

%% ============================================================
\section{Introduction}\label{sec:introduction}
%% ============================================================

The order in which training data is presented to a language model can
affect final performance. Curriculum learning approaches
\cite{bengio2009curriculum, graves2017automated} propose structured
orderings --- typically easy-to-hard --- to improve convergence or final
quality. Recent work has demonstrated measurable effects of data ordering
in LLM pretraining \cite{beyondrandom2025, lrdecay2025}, with
compression ratio and lexical diversity emerging as effective difficulty
signals.

Separately, framework-level compilation has become standard practice in
PyTorch training. \texttt{torch.compile} \cite{pytorch2compile} traces
execution graphs and generates fused CUDA kernels optimized for observed
data access patterns, typically yielding 1.5--2$\times$ throughput
improvements.

We report an unexpected interaction between these two developments:
the magnitude of data ordering effects depends strongly on whether
the training framework compiles execution graphs. Specifically, the
same curriculum experiment produces dramatic effects on
CUDA/\texttt{torch.compile} and null effects on Apple MLX, a framework
with no equivalent compilation step.

This finding emerged from an investigation of the King Wen I-Ching
sequence as a curriculum ordering \cite{chan2025kingwen}. Across 50+
controlled experiments on two hardware platforms, the curriculum
hypothesis was falsified --- but the platform discrepancy revealed
something unexpected about how compilation interacts with data structure.

%% ============================================================
\section{Experimental Setup}\label{sec:setup}
%% ============================================================

\subsection{Model and Data}

All experiments use the autoresearch framework \cite{karpathy2024autoresearch}:
a 4-layer GPT model ($\sim$3M parameters) with sliding window attention,
RoPE positional encoding, and a custom AdamW optimizer with per-parameter
learning rate groups. Training data is drawn from ClimbMix-400B via
parquet shards, tokenized with a BPE tokenizer (vocabulary size 8192)
and packed using best-fit document packing \cite{fewertruncations2024}.

Training is bounded by a fixed wall-clock budget of 300 seconds.
Curriculum overhead (difficulty scoring, buffer sorting) counts against
this budget, ensuring that expensive ordering strategies are
automatically penalized.

\subsection{Curriculum Buffer}

A 64-batch buffer wraps the dataloader. At each refill, all 64
micro-batches are scored by difficulty, reordered according to the
experimental condition, and yielded one at a time to the training loop.
Six ordering conditions are tested:

\begin{itemize}\tightlist
\item \textbf{Sequential} (no buffer): unmodified dataloader
\item \textbf{Buffered passthrough}: buffer enabled, no reordering
\item \textbf{Random shuffle}: buffer contents randomly permuted
\item \textbf{Easy-to-hard}: ascending difficulty
\item \textbf{Hard-to-easy}: descending difficulty
\item \textbf{King Wen}: rank-mapped to I-Ching surprise profile
\end{itemize}

\subsection{Difficulty Metrics}

Two independent metrics are tested to rule out metric-specific artifacts:

\begin{itemize}\tightlist
\item \textbf{Token diversity}: $\text{unique}(x) / |x|$ --- the fraction
of unique tokens in a micro-batch. Zero computational overhead.
\item \textbf{Compression ratio}: $|\text{gzip}(x.\text{bytes}())| / |x.\text{bytes}()|$
--- ratio of compressed to raw byte length. Approximately 7\% of
training budget as overhead.
\end{itemize}

\subsection{Platforms}

\begin{itemize}\tightlist
\item \textbf{CUDA}: Intel desktop, NVIDIA RTX 2060 (6 GB), PyTorch with
\texttt{torch.compile}, fp32 (Turing architecture lacks bf16 support)
\item \textbf{MLX}: MacBook Pro, Apple M-series, 96 GB unified memory,
MLX framework \cite{mlx2023}. No compilation step; operations execute
eagerly via Metal compute shaders.
\end{itemize}

Both platforms use identical model architecture, hyperparameters,
data shards, tokenizer, and evaluation procedure. The only differences
are the framework (PyTorch vs MLX) and hardware (NVIDIA vs Apple Silicon).

\subsection{Evaluation}

The primary metric is \textbf{val\_bpb} (validation bits per byte) ---
a vocabulary-size-independent measure of language modeling quality.
Lower is better. Evaluation uses $\sim$1.5M tokens from a held-out
validation shard. The seed noise floor is 0.04 bpb, established by
a 30-seed sweep (ADR-004 in supplementary materials).

%% ============================================================
\section{Results}\label{sec:results}
%% ============================================================

\subsection{CUDA: Strong Ordering Effects}

Table~\ref{tab:cuda} shows results on the CUDA platform with both
difficulty metrics.

\begin{table}[h]
\centering
\caption{CUDA results: val\_bpb by ordering and difficulty metric.
Seed 42, DEPTH=4, standard warmdown (0.5).}
\label{tab:cuda}
\begin{tabular}{lcccc}
\toprule
Ordering & Token Div & Comp Ratio & $\Delta$ (TD) & $\Delta$ (CR) \\
\midrule
Sequential       & 1.719 & 1.778 & ---    & ---    \\
Buffered pass    & 1.680 & 1.640 & -0.039 & -0.138 \\
\textbf{Random}  & \textbf{1.614} & \textbf{1.627} & \textbf{-0.106} & \textbf{-0.151} \\
Easy-to-hard     & 1.632 & 1.634 & -0.087 & -0.144 \\
Hard-to-easy     & 1.627 & 1.634 & -0.092 & -0.144 \\
King Wen         & 1.662 & 1.638 & -0.057 & -0.140 \\
\bottomrule
\end{tabular}
\end{table}

All orderings beat sequential. Random shuffle is the best ordering
under token diversity (0.106 improvement, 2.7$\times$ the noise floor).
Under compression ratio, the spread collapses to 0.011 but all orderings
still strongly beat sequential (0.138--0.151).

\subsection{MLX: Null Effects}

Table~\ref{tab:mlx} shows results on the MLX platform.

\begin{table}[h]
\centering
\caption{MLX results: val\_bpb by ordering and LR regime.
Seed 42, DEPTH=4, compression ratio metric.}
\label{tab:mlx}
\begin{tabular}{lcc}
\toprule
Ordering & Standard WD & Constant LR \\
\midrule
Sequential       & 1.732 & 1.722 \\
Random           & 1.713 & 1.697 \\
Easy-to-hard     & 1.695 & 1.731 \\
Hard-to-easy     & 1.709 & 1.707 \\
King Wen         & 1.724 & 1.729 \\
\bottomrule
\end{tabular}
\end{table}

No ordering beats sequential by more than the 0.060 noise floor
(measured via 3-seed sweep at DEPTH=4). The largest effect
(easy-to-hard: $-0.037$) is within noise. Token diversity on MLX
was also tested: sequential 1.663, random 1.673 --- random is
actually \emph{worse}, confirming the null result.

\subsection{Cross-Platform Comparison}

\begin{table}[h]
\centering
\caption{Cross-platform comparison (compression ratio metric).
$\Delta$ is improvement over each platform's sequential baseline.}
\label{tab:crossplatform}
\begin{tabular}{lcccr}
\toprule
Ordering & CUDA & MLX & $\Delta$ CUDA & $\Delta$ MLX \\
\midrule
Sequential   & 1.778 & 1.732 & ---    & ---    \\
Random       & 1.627 & 1.713 & -0.151 & -0.020 \\
Easy-to-hard & 1.634 & 1.695 & -0.144 & -0.037 \\
King Wen     & 1.638 & 1.724 & -0.140 & -0.009 \\
\bottomrule
\end{tabular}
\end{table}

The platform gap is 4--10$\times$ larger than the metric gap.
CUDA shows 0.138--0.151 bpb improvement from any reordering.
MLX shows 0.009--0.037 --- within seed noise.
The effect is platform-dependent, not metric-dependent.

%% ============================================================
\section{Analysis}\label{sec:analysis}
%% ============================================================

\subsection{Ruling Out Confounds}

Several alternative explanations are considered and ruled out:

\textbf{Hardware differences:} The CUDA machine (RTX 2060, 6 GB) is
weaker than the MLX machine (96 GB unified memory). If anything, the
more powerful machine should show larger effects. It shows smaller ones.

\textbf{Precision:} The CUDA machine trains in fp32 (Turing lacks bf16).
The MLX machine also uses fp32. Precision is controlled.

\textbf{Difficulty metric:} Both token diversity and compression ratio
were tested on both platforms. The platform gap persists across both
metrics.

\textbf{Evaluation noise:} The MLX noise floor (0.060, measured via
3-seed sweep) is larger than CUDA's (0.040, 30-seed sweep), but the
CUDA effects (0.106--0.151) are 2.5--3.8$\times$ their noise floor,
while MLX effects (0.009--0.037) are 0.2--0.6$\times$.

\textbf{Data ordering:} Both platforms use the same dataloader, the
same best-fit packing, and the same buffer implementation. The data
is identically ordered on both machines.

\subsection{The Compilation Hypothesis}

The remaining systematic difference is \texttt{torch.compile}.
We propose that the compilation process creates the observed interaction:

\begin{enumerate}
\item Best-fit document packing creates implicit sequential correlation
in the training data: adjacent batches share similar document-length
distributions and content patterns \cite{fewertruncations2024}.

\item \texttt{torch.compile} traces execution graphs and generates
fused CUDA kernels optimized for the data access patterns it observes.
When data has consistent sequential structure, the compiled kernels
specialize to that structure.

\item This specialization creates a form of \emph{kernel-level
overfitting}: the compiled computation graph is locally optimal for the
sequential data pattern but suboptimal for the general distribution.

\item Shuffling the buffer disrupts the sequential pattern. The compiled
kernels, now mismatched to the data, either recompile or fall back to
more general execution paths. These general paths, while individually
slower, produce better-generalizing gradient updates.
\end{enumerate}

On MLX, there is no compilation step. Operations execute eagerly via
Metal compute shaders that do not adapt to data access patterns.
Shuffling has no mechanism to disrupt because there is no
pattern-adapted computation to disrupt.

\subsection{Connection to Memorization-Compression Dynamics}

Recent work on memorization-compression cycles \cite{gapt2025} has shown
that LLM training naturally alternates between memorization
(representation expansion) and compression (entropy reduction) phases.
The gradient alignment between cross-entropy and representation entropy
oscillates throughout training.

The compilation hypothesis suggests a mechanism by which
\texttt{torch.compile} may amplify the memorization phase. Compiled
kernels that specialize to sequential patterns effectively ``memorize''
the data structure at the computation level, reducing the implicit
compression pressure that comes from processing diverse data.
Shuffling restores this compression pressure by forcing the computation
graph to handle unpredictable access patterns.

This interpretation aligns with the observation that buffered passthrough
alone (without scoring or reordering) provides substantial benefit on CUDA
(0.039--0.138 bpb). The mere act of disrupting sequential access patterns
--- even without any curriculum logic --- accounts for most of the effect.

%% ============================================================
\section{Implications}\label{sec:implications}
%% ============================================================

\subsection{For Curriculum Learning Research}

Any curriculum learning experiment conducted with \texttt{torch.compile}
(or similar compilation frameworks such as JAX/XLA \cite{jax2018})
should include a non-compiled baseline to distinguish genuine curriculum
effects from compilation artifacts. Studies reporting curriculum
improvements without controlling for compilation may be measuring
an interaction effect rather than a learning dynamics effect.

\subsection{For Training Infrastructure}

The finding that shuffling improves performance on compiled frameworks
suggests a practical optimization: inserting a random shuffle buffer
between the dataloader and the training loop may yield free performance
gains on compiled pipelines, independent of any curriculum design.
This is particularly relevant for best-fit packed dataloaders, which
create the sequential correlation that compilation exploits.

\subsection{For Framework Developers}

The kernel-level overfitting mechanism suggests that compilation
strategies could be modified to avoid specializing to data-dependent
access patterns. Periodic kernel recompilation or input-dependent
dispatch could mitigate the effect without requiring explicit shuffling.

%% ============================================================
\section{Limitations}\label{sec:limitations}
%% ============================================================

\textbf{Scale:} All experiments use a 3M-parameter model with a 5-minute
training budget. The compilation hypothesis may not hold at larger scales
where kernel specialization is a smaller fraction of total computation.

\textbf{Mechanism verification:} The compilation hypothesis is
consistent with observations but has not been verified at the kernel
level. Direct analysis of compiled CUDA kernels under sequential vs
shuffled data would strengthen the claim.

\textbf{Other compiled frameworks:} Only PyTorch \texttt{torch.compile}
and MLX are compared. Testing on JAX/XLA, TensorFlow XLA, or
TorchScript would establish generality.

\textbf{Single model architecture:} Only one GPT variant is tested.
The interaction may differ for models with different attention patterns,
batch normalization, or non-standard layer structures.

%% ============================================================
\section{Related Work}\label{sec:related}
%% ============================================================

\textbf{Curriculum learning} has been studied extensively in LLM
pretraining. \citet{beyondrandom2025} demonstrate up to 3.5\% gains with
compression-ratio difficulty. \citet{lrdecay2025} show that standard LR
decay suppresses curriculum benefits by up to 44$\times$.
\citet{influence2025} propose model-centric difficulty scoring.
\citet{yang2025estimating} estimate ordering effects without retraining.

\textbf{Compilation effects on training} have received limited attention.
\citet{pytorch2compile} describe \texttt{torch.compile}'s architecture
but focus on throughput rather than training dynamics.
To our knowledge, no prior work documents an interaction between
compilation and data ordering effects.

\textbf{Data packing and ordering artifacts.}
\citet{fewertruncations2024} document that best-fit packing creates
deterministic ordering biases but do not investigate shuffling as a
remedy. Our work identifies shuffling as effective specifically because
it disrupts compilation-level specialization, not because it improves
the curriculum per se.

\textbf{Memorization-compression dynamics.}
\citet{gapt2025} show that training alternates between memorization and
compression phases, with oscillating gradient alignment. Our compilation
hypothesis suggests that kernel specialization may amplify the
memorization phase on compiled frameworks.

%% ============================================================
\section{Conclusion}\label{sec:conclusion}
%% ============================================================

We document a previously unreported interaction between
\texttt{torch.compile} and training data ordering.
Curriculum-like effects that appear large (0.106--0.151 bpb) on
compiled CUDA training are absent (0.009--0.037 bpb, within noise)
on non-compiled MLX training, despite identical model architecture,
data, and evaluation.

The most parsimonious explanation is that \texttt{torch.compile}'s
kernel optimization specializes to sequential data access patterns
created by best-fit packing. Any disruption of these patterns ---
random shuffling, structured reordering, or even buffered passthrough
--- forces more general computation that produces better-generalizing
gradient updates.

This finding has immediate practical implications: a random shuffle
buffer provides free performance gains on compiled pipelines, and
curriculum learning experiments should control for compilation to avoid
reporting spurious effects.

\section*{Acknowledgments}

This work emerged from the King Wen I-Ching experimental research
program \cite{chan2025kingwen}. The author thanks Claude (Anthropic)
for assistance with experimental design and analysis. All code and
data are open source at
\url{https://github.com/augchan42/autoresearch} (CUDA) and
\url{https://github.com/augchan42/autoresearch-mlx} (MLX).

\bibliographystyle{plain}
\bibliography{references}

\end{document}
